\section{Numerical results}
\label{sec:num}

In this section we numerically analyze the quasi-PDF, \spcorr and pseudo-PDF by studying how the matching coefficients in Eqs.~(\ref{eq:fac-tq},\ref{eq:ps-q-fact}) change the PDF. The quasi-PDF has already been studied in this manner for the $\overline{\rm MS}$, transverse momentum cut-off, and RI/MOM schemes in Ref.~\cite{Stewart:2017tvs}. Our  new $\overline{\rm MS}$ result for the matching is given in \eq{quasi-matching}, and leads to stable convolution integrals.  We also compare the differences between using hadron momentum $p^z=P^z$ and the parton momentum $p^z=|y|P^z$ for the matching coefficient in the $\overline{\rm MS}$ scheme. We take $\Gamma =\gamma^0$ for the results here.

As an example we use for our analysis the unpolarized iso-vector parton distribution,
\begin{align}
f_{u-d}(x,\mu) =  f_u(x,\mu) - f_d(x,\mu) - f_{\bar{u}}(-x,\mu) +  f_{\bar{d}}(-x,\mu),
\end{align}
where we include
$f_{\bar{u}}(-x,\mu) = - f_{\bar{u}}(x,\mu)$ and $f_{\bar{d}}(-x,\mu) = - f_{\bar{d}}(x,\mu)$, the anti-parton distributions. For ease of comparison, we use the next-to-leading-order iso-vector PDF $f_{u-d}$ from MSTW 2008~\cite{Martin:2009iq} with the corresponding running coupling $\alpha_s(\mu)$.

To implement the plus functions in the numerical calculation, we impose a soft cutoff $|y-x|<10^{-m}$ and test the sensitivity of results to $m$. Since the limit of $y\to0$ corresponds to the asymptotic region $|x/y|\to\infty$, we also impose a UV cutoff $|y|>10^{-n}$ to test the convergence of the convolution integral. We find that all the results presented below are insensitive to $m$ and $n$. The fact that our result in \eq{quasi-matching} has terms outside the plus function at $1$ in each of the $\xi\in [1,\infty]$ and $\xi\in [-\infty,0]$ intervals is important for ensuring that our $\overline{\rm MS}$ result for $C$ is insensitive to the $|y|>10^{-n}$ cutoff. This was not the case for the quasi-PDF that was defined with a transverse momentum cutoff~\cite{Xiong:2013bka}. The RI/MOM scheme result~\cite{Stewart:2017tvs} also does not suffer from this issue.

In Fig.~\ref{fig:quasi} we compare the PDF with the quasi-PDF in the $\overline{\rm MS}$ scheme obtained from the convolution in \eq{fac-tq} using our one-loop result in \eq{quasi-matching}. We observe that changing from $p^z=P^z$  to the correct $p^z=|y|P^z$ shifts the result in the physical region by a considerable amount.

\begin{figure*}
	\centering
	\includegraphics[width=.49\textwidth]{images/Compare_iPDFvsPDFRe_TeX.pdf}
	\hspace{0.1cm}
	\includegraphics[width=0.49\textwidth]{images/Compare_iPDFvsPDFRe_z_TeX.pdf}
	\\
	\includegraphics[width=.49\textwidth]{images/Compare_iPDFvsPDFIm_TeX.pdf}
	\hspace{0.1cm}
	\includegraphics[width=0.49\textwidth]{images/Compare_iPDFvsPDFIm_z_TeX.pdf}	\\
\vspace{-0.3cm}
	\caption{(left) Comparison between the light-cone time distribution $Q_{u-d}$ and \spcorr from $({\cal C}^{\overline{\rm MS}}\otimes Q_{u-d})$ in the $\overline{\rm MS}$ scheme. The orange and blue bands indicate the results from varying the factorization scale $\mu=4\,{\rm GeV}$ by a factor of two. (right) Same but now showing only central \spcorr curves for different values of $z$. The top panels show the real part, while bottom panels show the imaginary part. }
	\label{fig:ioffe}
\end{figure*}

The same type of comparison can be made for the pseudo-PDF in the $\overline{\rm MS}$ scheme by applying the factorization formula in Eq.~(\ref{eq:ps-q-fact}) and matching coefficient in Eq.~(\ref{eq:ps-c}). In Fig.~\ref{fig:pseudo}a we compare the PDF and pseudo-PDF and their dependence on the factorization scale $\mu$, while in Fig.~\ref{fig:pseudo}b we include the dependence of the pseudo-PDF on the distance $|z|$. Since the matching coefficient in Eq.(\ref{eq:ps-c}) is similar to the parton splitting function except for the nontrivial finite constants, matching the PDF to the pseudo-PDF is analogous to evolving the PDF from $\mu$ to the scale of $1/|z|$. This evolution has been calculated in Refs.~\cite{Radyushkin:2017cyf}.  The variation of $|z|$ has a similar effect to the PDF evolution, as is observed in the right panel of Fig.~\ref{fig:pseudo}. When $|z|\mu=1$, the logarithm is zero, and the matching effect from ${\cal C}$ is determined by the nontrivial constants in Eq.(\ref{eq:ps-c}), which shifts the PDF downward in the large-$x$ region and upward in the small-$x$ region.

Finally, we can make a similar comparison for the \spcorr in the $\overline{\rm MS}$ scheme obtained with Eq.~(\ref{eq:io-Q-fact}) and Eq.~(\ref{eq:ps-c}). Its real and imaginary parts are even and odd functions of $\zeta$ respectively, and are shown in Fig.~\ref{fig:ioffe}. Again we show the residual dependence on $\mu$ and $|z|$ which are similar to that for the pseudo-PDF.   The matching broadens the curves in the coordinate space.  The \spcorr renormalized in the $\overline{\rm MS}$ scheme does not exhibit vector current (or particle number) conservation, which can be clearly seen from the fact that the real part of the distribution is not equal to 1 at $\zeta=0$ (except for the special case where $|z|\mu$ is tuned to cancel the constant terms in the one-loop ${\cal C}$).
